\documentclass[11pt,a4paper]{article}
\usepackage[utf8]{inputenc}
\usepackage[T1]{fontenc}
\usepackage[french]{babel}
\usepackage{amsmath,amssymb,amsfonts,amsthm,mathrsfs,mathtools}
\usepackage[top=2cm,bottom=2cm,left=2cm,right=2cm]{geometry}
\usepackage[dvipsnames]{xcolor}
\usepackage{tcolorbox}
\tcbuselibrary{skins,breakable}
\usepackage{tikz}
\usetikzlibrary{calc,shapes,arrows.meta}
\usepackage{enumitem}
\usepackage{fancyhdr}
\usepackage{lastpage}
\usepackage{hyperref}

% Couleurs personnalisées sécurisées
\definecolor{emeraldframe}{RGB}{20, 110, 75}
\definecolor{emeraldback}{RGB}{242, 250, 246}
\definecolor{emeraldtitle}{RGB}{15, 85, 55}
\definecolor{prepaframe}{RGB}{30, 70, 130}
\definecolor{prepaback}{RGB}{245, 248, 255}
\definecolor{prepatitle}{RGB}{20, 50, 100}

% Configuration des hyperliens
\hypersetup{
    colorlinks=true,
    linkcolor=blue!70!black,
    urlcolor=blue!70!black,
    citecolor=blue!70!black
}

% En-têtes et pieds de page
\pagestyle{fancy}
\fancyhf{}
\renewcommand{\headrulewidth}{0.5pt}
\renewcommand{\footrulewidth}{0.4pt}
\fancyhead[L]{\small\textbf{Lycée Louis-le-Grand} -- 1\textsuperscript{ère} Spécialité (1E3)}
\fancyhead[C]{\small\textbf{DST 3 : Équations de droites}}
\fancyhead[R]{\small 25 Novembre 2021}
\fancyfoot[L]{\footnotesize LLG 2021/2022}
\fancyfoot[C]{\footnotesize Énoncés \& Corrigés ultra-détaillés}
\fancyfoot[R]{\footnotesize Page \thepage\ / \pageref{LastPage}}

% Boîtes tcolorbox
\newtcolorbox{correctionbox}[1][]{
  colback=emeraldback,
  colframe=emeraldframe,
  fonttitle=\bfseries\color{white},
  coltitle=white,
  colbacktitle=emeraldtitle,
  title={Correction détaillée et méthodologie #1},
  arc=2mm,
  boxrule=0.8pt,
  titlerule=0.5pt,
  breakable,
  before skip=0.4cm,
  after skip=0.5cm
}

\newtcolorbox{remarquebox}[1][]{
  colback=prepaback,
  colframe=prepaframe,
  fonttitle=\bfseries\small\color{white},
  coltitle=white,
  colbacktitle=prepatitle,
  title={Prolongement \& Regard Prépa / MPSI #1},
  arc=1.5mm,
  boxrule=0.6pt,
  titlerule=0.4pt,
  breakable,
  before skip=0.2cm,
  after skip=0.3cm
}

% Commandes utiles
\newcommand{\N}{\mathbb{N}}
\newcommand{\R}{\mathbb{R}}
\newcommand{\vect}[1]{\overrightarrow{#1}}
\newcommand{\dd}{\mathrm{d}}

\begin{document}

% Titre principal
\begin{center}
    {\LARGE \bfseries Devoir Sur Table n°3 : Équations de droites}\\[0.25cm]
    {\large \textbf{Classe de Première Spécialité Mathématiques -- Classe 1E3}}\\[0.15cm]
    {\small Lycée Louis-le-Grand -- Jeudi 25 Novembre 2021 -- Durée : 2 heures}\\[0.3cm]
    \hrule height 1pt
\end{center}

\vspace{0.3cm}

\begin{tcolorbox}[colback=gray!5!white,colframe=gray!60!black,arc=1.5mm,boxrule=0.6pt]
\textbf{Consignes de rédaction :}
\begin{itemize}[leftmargin=*,itemsep=1pt]
    \item La clarté, la rigueur des justifications et le soin apporté à la rédaction constituent une part essentielle de l'évaluation.
    \item Tout résultat numérique ou géométrique doit être rigoureusement démontré en s'appuyant sur les théorèmes du cours (colinéarité, caractérisation par déterminant, système d'équations).
\end{itemize}
\end{tcolorbox}

\vspace{0.3cm}

% ====================================================
% EXERCICE 1
% ====================================================
\subsection*{Exercice 1 \hfill \normalfont\normalsize (4 points)}

On se place dans un repère du plan.

On considère les points $A(-7\,;\,3)$, $B(8\,;\,-3)$, $C(3\,;\,-1)$ et $D(-2\,;\,1)$.

\begin{enumerate}
    \item Établir une équation cartésienne de la droite $(AB)$ sans calcul de coefficient directeur.
    \item 
    \begin{enumerate}
        \item Calculer le coefficient directeur de la droite $(CD)$.
        \item En déduire l'ordonnée à l'origine de la droite $(CD)$ puis donner son équation réduite.
    \end{enumerate}
    \item Montrer que les droites $(AB)$ et $(CD)$ sont confondues.
\end{enumerate}

\begin{correctionbox}
\begin{enumerate}
    \item \textbf{Équation cartésienne de la droite $(AB)$ par le déterminant :}
    
    Soit $M(x\,;\,y)$ un point quelconque du plan.
    
    Les coordonnées du vecteur $\vect{AB}$ sont :
    \[
    \vect{AB} = \begin{pmatrix} x_B - x_A \\ y_B - y_A \end{pmatrix} = \begin{pmatrix} 8 - (-7) \\ -3 - 3 \end{pmatrix} = \begin{pmatrix} 15 \\ -6 \end{pmatrix}.
    \]
    Le vecteur $\vect{AB}$ n'est pas le vecteur nul (car $15 \neq 0$), donc la droite $(AB)$ est bien définie et admet $\vect{AB}$ comme vecteur directeur.
    
    Les coordonnées du vecteur $\vect{AM}$ sont :
    \[
    \vect{AM} = \begin{pmatrix} x - x_A \\ y - y_A \end{pmatrix} = \begin{pmatrix} x - (-7) \\ y - 3 \end{pmatrix} = \begin{pmatrix} x + 7 \\ y - 3 \end{pmatrix}.
    \]
    Par définition, $M(x\,;\,y) \in (AB)$ si et seulement si les vecteurs $\vect{AM}$ et $\vect{AB}$ sont colinéaires, ce qui équivaut à la nullité de leur déterminant :
    \[
    M(x\,;\,y) \in (AB) \iff \det(\vect{AM}, \vect{AB}) = 0.
    \]
    Or :
    \[
    \det(\vect{AM}, \vect{AB}) = \begin{vmatrix} x + 7 & 15 \\ y - 3 & -6 \end{vmatrix} = -6(x + 7) - 15(y - 3).
    \]
    Développons et réduisons cette expression :
    \[
    -6(x + 7) - 15(y - 3) = -6x - 42 - 15y + 45 = -6x - 15y + 3.
    \]
    Ainsi :
    \[
    M(x\,;\,y) \in (AB) \iff -6x - 15y + 3 = 0.
    \]
    En divisant cette équation par $-3$, on obtient une forme irréductible équivalente :
    \[
    \mathbf{2x + 5y - 1 = 0}.
    \]
    Une équation cartésienne de la droite $(AB)$ est donc :
    \[
    (AB) : 2x + 5y - 1 = 0 \quad (\text{ou } -6x - 15y + 3 = 0).
    \]

    \item 
    \begin{enumerate}
        \item \textbf{Coefficient directeur de $(CD)$ :}
        
        Les points $C(3\,;\,-1)$ et $D(-2\,;\,1)$ vérifient $x_C \neq x_D$ ($3 \neq -2$). La droite $(CD)$ n'est donc pas parallèle à l'axe des ordonnées.
        
        Son coefficient directeur $m$ est donné par :
        \[
        m = \frac{y_D - y_C}{x_D - x_C} = \frac{1 - (-1)}{-2 - 3} = \frac{2}{-5} = -\frac{2}{5}.
        \]

        \item \textbf{Ordonnée à l'origine et équation réduite de $(CD)$ :}
        
        L'équation réduite de $(CD)$ est de la forme $y = mx + p$, c'est-à-dire :
        \[
        y = -\frac{2}{5}x + p, \quad \text{avec } p \in \R.
        \]
        Puisque $C(3\,;\,-1) \in (CD)$, ses coordonnées vérifient cette relation :
        \[
        -1 = -\frac{2}{5} \times 3 + p \iff -1 = -\frac{6}{5} + p \iff p = -1 + \frac{6}{5} = \frac{1}{5}.
        \]
        \textit{Vérification avec le point $D(-2\,;\,1)$ :}
        \[
        -\frac{2}{5} \times (-2) + \frac{1}{5} = \frac{4}{5} + \frac{1}{5} = \frac{5}{5} = 1 = y_D. \quad \text{(Vérifié)}
        \]
        L'ordonnée à l'origine de $(CD)$ est donc $p = \dfrac{1}{5}$, et son équation réduite est :
        \[
        (CD) : y = -\frac{2}{5}x + \frac{1}{5}.
        \]
    \end{enumerate}

    \item \textbf{Droites confondues :}
    
    Déterminons l'équation réduite de la droite $(AB)$ à partir de son équation cartésienne :
    \[
    2x + 5y - 1 = 0 \iff 5y = -2x + 1 \iff y = -\frac{2}{5}x + \frac{1}{5}.
    \]
    Les droites $(AB)$ et $(CD)$ possèdent exactement la même équation réduite $y = -\dfrac{2}{5}x + \dfrac{1}{5}$.
    
    Par conséquent, l'ensemble des points de $(AB)$ coïncide avec l'ensemble des points de $(CD)$ :
    \[
    \mathbf{(AB) \text{ et } (CD) \text{ sont confondues}}.
    \]
\end{enumerate}
\end{correctionbox}

\vspace{0.3cm}

% ====================================================
% EXERCICE 2
% ====================================================
\subsection*{Exercice 2 \hfill \normalfont\normalsize (4{,}5 points)}

\begin{enumerate}
    \item Donner la hauteur d'un triangle équilatéral de côté 1.
    \item On considère la configuration géométrique suivante :
    
    \begin{center}
    \begin{tikzpicture}[scale=3.2, cap=round, >=Latex]
        % Sommets du carré
        \coordinate (A) at (0,0);
        \coordinate (B) at (1,0);
        \coordinate (C) at (1,1);
        \coordinate (D) at (0,1);
        
        % Sommet E (intérieur) : hauteur h = sqrt(3)/2
        \coordinate (E) at (0.5, {sqrt(3)/2});
        
        % Sommet F (extérieur) : x = 1 + sqrt(3)/2, y = 0.5
        \coordinate (F) at ({1 + sqrt(3)/2}, 0.5);
        
        % Tracé du carré ABCD
        \draw[thick] (A) -- (B) -- (C) -- (D) -- cycle;
        
        % Tracé triangle équilatéral AEB
        \draw[thick] (A) -- (E) -- (B);
        
        % Tracé triangle équilatéral BCF
        \draw[thick] (B) -- (F) -- (C);
        
        % Droite passant par D, E, F (en pointillés rouge pour illustrer l'alignement)
        \draw[dashed, red!80!black, thick] ($(D)!-0.1!(F)$) -- ($(D)!1.1!(F)$);
        
        % Repère : vecteurs AB et AD
        \draw[->, very thick, blue!80!black] (A) -- (0.5, 0) node[below] {\small $\vect{AB}$};
        \draw[->, very thick, blue!80!black] (A) -- (0, 0.5) node[left] {\small $\vect{AD}$};
        
        % Points et labels
        \fill (A) circle (0.6pt) node[below left] {$A$};
        \fill (B) circle (0.6pt) node[below right] {$B$};
        \fill (C) circle (0.6pt) node[above right] {$C$};
        \fill (D) circle (0.6pt) node[above left] {$D$};
        \fill (E) circle (0.6pt) node[below] {$E$};
        \fill (F) circle (0.6pt) node[right] {$F$};
    \end{tikzpicture}
    \end{center}
    
    $ABCD$ est un carré. $AEB$ et $BCF$ sont des triangles équilatéraux.
    \begin{enumerate}
        \item Donner les coordonnées de $D$, $E$ et $F$ dans un repère bien choisi.
        \item En déduire que les points $D$, $E$ et $F$ sont alignés.
    \end{enumerate}
\end{enumerate}

\begin{correctionbox}
\begin{enumerate}
    \item \textbf{Hauteur d'un triangle équilatéral de côté 1 :}
    
    Soit un triangle équilatéral de côté $a = 1$. La hauteur issue d'un sommet est également médiatrice du côté opposé, donc elle coupe ce côté en son milieu.
    
    Cette hauteur découpe ainsi le triangle en deux triangles rectangles dont l'hypoténuse mesure 1 et le côté adjacent mesure $\dfrac{1}{2}$.
    
    D'après le théorème de Pythagore, la hauteur $h$ vérifie :
    \[
    h^2 + \left(\frac{1}{2}\right)^2 = 1^2 \iff h^2 = 1 - \frac{1}{4} = \frac{3}{4}.
    \]
    Comme $h > 0$, on a :
    \[
    \mathbf{h = \frac{\sqrt{3}}{2}}.
    \]

    \item 
    \begin{enumerate}
        \item \textbf{Choix du repère et détermination des coordonnées :}
        
        Puisque $ABCD$ est un carré, on a $(AB) \perp (AD)$ et $AB = AD = 1$. Le repère $\mathcal{R} = (A\,;\, \vect{AB}, \vect{AD})$ est donc un \textbf{repère orthonormé}.
        
        Dans ce repère :
        \begin{itemize}[leftmargin=*,itemsep=2pt]
            \item $A$ est l'origine du repère : $A(0\,;\,0)$.
            \item $\vect{AB} = 1\vect{AB} + 0\vect{AD}$, donc $B(1\,;\,0)$.
            \item $\vect{AD} = 0\vect{AB} + 1\vect{AD}$, donc $\mathbf{D(0\,;\,1)}$.
            \item $\vect{AC} = \vect{AB} + \vect{AD}$, donc $C(1\,;\,1)$.
            \item \textbf{Point $E$ :} Le triangle $AEB$ est équilatéral de base $[AB]$.
            Le projeté orthogonal de $E$ sur l'axe des abscisses $(AB)$ est le milieu de $[AB]$, donc $x_E = \dfrac{1}{2}$.
            L'ordonnée de $E$ correspond à la hauteur du triangle $AEB$, dirigée vers le haut, donc $y_E = h = \dfrac{\sqrt{3}}{2}$.
            D'où :
            \[
            \mathbf{E\left(\frac{1}{2}\,;\, \frac{\sqrt{3}}{2}\right)}.
            \]
            \item \textbf{Point $F$ :} Le triangle $BCF$ est équilatéral de base $[BC]$.
            Le projeté orthogonal de $F$ sur $(BC)$ est le milieu de $[BC]$, donc son ordonnée est $y_F = \dfrac{1}{2}$.
            Son abscisse est égale à l'abscisse de la droite $(BC)$ (qui vaut 1) augmentée de la hauteur $h$ du triangle $BCF$ :
            $x_F = 1 + h = 1 + \dfrac{\sqrt{3}}{2}$.
            D'où :
            \[
            \mathbf{F\left(1 + \frac{\sqrt{3}}{2}\,;\, \frac{1}{2}\right)}.
            \]
        \end{itemize}

        \item \textbf{Alignement des points $D$, $E$ et $F$ :}
        
        Calculons les coordonnées des vecteurs $\vect{DE}$ et $\vect{DF}$ :
        \[
        \vect{DE} = \begin{pmatrix} x_E - x_D \\ y_E - y_D \end{pmatrix} = \begin{pmatrix} \dfrac{1}{2} - 0 \\ \dfrac{\sqrt{3}}{2} - 1 \end{pmatrix} = \begin{pmatrix} \dfrac{1}{2} \\ \dfrac{\sqrt{3}}{2} - 1 \end{pmatrix},
        \]
        \[
        \vect{DF} = \begin{pmatrix} x_F - x_D \\ y_F - y_D \end{pmatrix} = \begin{pmatrix} 1 + \dfrac{\sqrt{3}}{2} - 0 \\ \dfrac{1}{2} - 1 \end{pmatrix} = \begin{pmatrix} 1 + \dfrac{\sqrt{3}}{2} \\ -\dfrac{1}{2} \end{pmatrix}.
        \]
        Calculons le déterminant $\det(\vect{DE}, \vect{DF})$ :
        \[
        \det(\vect{DE}, \vect{DF}) = \begin{vmatrix} \dfrac{1}{2} & 1 + \dfrac{\sqrt{3}}{2} \\ \dfrac{\sqrt{3}}{2} - 1 & -\dfrac{1}{2} \end{vmatrix} = \frac{1}{2} \times \left(-\frac{1}{2}\right) - \left(\frac{\sqrt{3}}{2} - 1\right)\left(1 + \frac{\sqrt{3}}{2}\right).
        \]
        Or, en reconnaissant l'identité remarquable $(a - b)(a + b) = a^2 - b^2$ avec $a = \dfrac{\sqrt{3}}{2}$ et $b = 1$ :
        \[
        \left(\frac{\sqrt{3}}{2} - 1\right)\left(\frac{\sqrt{3}}{2} + 1\right) = \left(\frac{\sqrt{3}}{2}\right)^2 - 1^2 = \frac{3}{4} - 1 = -\frac{1}{4}.
        \]
        Par conséquent :
        \[
        \det(\vect{DE}, \vect{DF}) = -\frac{1}{4} - \left(-\frac{1}{4}\right) = -\frac{1}{4} + \frac{1}{4} = 0.
        \]
        Puisque $\det(\vect{DE}, \vect{DF}) = 0$, les vecteurs $\vect{DE}$ et $\vect{DF}$ sont colinéaires.
        
        Ayant le point $D$ en commun, on conclut que :
        \[
        \mathbf{\text{Les points } D, E \text{ et } F \text{ sont alignés}}.
        \]
    \end{enumerate}
\end{enumerate}
\end{correctionbox}

\begin{remarquebox}[Démonstration géométrique synthétique classique]
En géométrie euclidienne pure sans coordonnées, ce résultat se démontre par le calcul des angles :
\begin{itemize}[leftmargin=*,itemsep=1pt]
    \item Dans le triangle isocèle $ADE$ en $A$ : $\widehat{DAE} = \widehat{DAB} - \widehat{EAB} = 90^\circ - 60^\circ = 30^\circ$.
    Comme $AD = AE = 1$, le triangle $ADE$ est isocèle en $A$, d'où $\widehat{ADE} = \dfrac{180^\circ - 30^\circ}{2} = 75^\circ$.
    \item Dans le triangle équilatéral $BCF$ et le carré $ABCD$ : $\widehat{DCF} = \widehat{DCB} + \widehat{BCF} = 90^\circ + 60^\circ = 150^\circ$.
    Comme $CD = CF = 1$, le triangle $DCF$ est isocèle en $C$, d'où $\widehat{CDF} = \dfrac{180^\circ - 150^\circ}{2} = 15^\circ$.
    \item Par conséquent, $\widehat{EDF} = \widehat{ADE} + \widehat{EDC} \dots$ ou plus simplement :
    $\widehat{ADE} + \widehat{CDF} = 75^\circ + 15^\circ = 90^\circ = \widehat{ADC}$.
    Les demi-droites $[DE)$ et $[DF)$ forment donc le même angle avec $(DC)$, ce qui prouve immédiatement l'alignement des points $D$, $E$ et $F$.
\end{itemize}
\end{remarquebox}

\vspace{0.3cm}

% ====================================================
% EXERCICE 3
% ====================================================
\subsection*{Exercice 3 \hfill \normalfont\normalsize (4{,}5 points)}

Soit un parallélogramme $ABCD$. On définit les points $E, F$ et $G$ par :
\[
\vect{DE} = 2\vect{AD}, \quad 3\vect{FD} + \vect{FC} = \vect{0} \quad \text{et} \quad \vect{AG} = m\vect{AF} \quad (m \in \R).
\]
On se place dans le repère $(A\,;\, \vect{AB}, \vect{AD})$.

\begin{enumerate}
    \item Déterminer les coordonnées de $A, B, C$ et $D$ sans justification.
    \item Calculer les coordonnées des points $E$ et $F$.
    \item Exprimer les coordonnées du point $G$ en fonction de $m$.
    \item Déterminer la valeur de $m$ pour que les points $E, G$ et $B$ soient alignés.
\end{enumerate}

\begin{correctionbox}
\begin{enumerate}
    \item \textbf{Coordonnées directes des sommets du parallélogramme :}
    
    Dans le repère $(A\,;\, \vect{AB}, \vect{AD})$ :
    \[
    \mathbf{A(0\,;\,0)}, \quad \mathbf{B(1\,;\,0)}, \quad \mathbf{C(1\,;\,1)}, \quad \mathbf{D(0\,;\,1)}.
    \]
    \textit{(En effet, par définition d'un parallélogramme, $\vect{AC} = \vect{AB} + \vect{AD} = 1\vect{AB} + 1\vect{AD}$)}.

    \item \textbf{Coordonnées des points $E$ et $F$ :}
    
    \begin{itemize}[leftmargin=*,itemsep=3pt]
        \item \textbf{Pour le point $E$ :}
        On a la relation vectorielle $\vect{DE} = 2\vect{AD}$.
        
        En passant aux coordonnées :
        \[
        \begin{cases}
        x_E - x_D = 2(x_D - x_A) \\
        y_E - y_D = 2(y_D - y_A)
        \end{cases}
        \iff
        \begin{cases}
        x_E - 0 = 2(0 - 0) = 0 \\
        y_E - 1 = 2(1 - 0) = 2
        \end{cases}
        \iff
        \begin{cases}
        x_E = 0 \\
        y_E = 3
        \end{cases}
        \]
        D'où :
        \[
        \mathbf{E(0\,;\,3)}.
        \]
        \textit{(Remarque vectorielle directe : $\vect{AE} = \vect{AD} + \vect{DE} = \vect{AD} + 2\vect{AD} = 3\vect{AD}$, ce qui donne immédiatement $(0\,;\,3)$)}.

        \item \textbf{Pour le point $F$ :}
        On a la relation vectorielle $3\vect{FD} + \vect{FC} = \vect{0}$.
        
        En coordonnées :
        \[
        \begin{cases}
        3(x_D - x_F) + (x_C - x_F) = 0 \\
        3(y_D - y_F) + (y_C - y_F) = 0
        \end{cases}
        \iff
        \begin{cases}
        3(0 - x_F) + (1 - x_F) = 0 \\
        3(1 - y_F) + (1 - y_F) = 0
        \end{cases}
        \]
        Résolvons chacune des deux équations scalaires :
        \[
        -3x_F + 1 - x_F = 0 \iff -4x_F + 1 = 0 \iff 4x_F = 1 \iff x_F = \frac{1}{4}.
        \]
        \[
        3 - 3y_F + 1 - y_F = 0 \iff 4 - 4y_F = 0 \iff 4y_F = 4 \iff y_F = 1.
        \]
        D'où :
        \[
        \mathbf{F\left(\frac{1}{4}\,;\,1\right)}.
        \]
        \textit{(Remarque barycentrique : $F$ est le barycentre de $\{(D, 3), (C, 1)\}$, donc $\vect{DF} = \dfrac{1}{4}\vect{DC} = \dfrac{1}{4}\vect{AB}$)}.
    \end{itemize}

    \item \textbf{Coordonnées du point $G$ en fonction de $m$ :}
    
    On a la relation vectorielle $\vect{AG} = m\vect{AF}$.
    
    Comme $A(0\,;\,0)$ est l'origine du repère, les coordonnées du point $G$ sont données par les composantes du vecteur $\vect{AG}$ :
    \[
    \begin{cases}
    x_G = m\, x_F = m \times \dfrac{1}{4} = \dfrac{m}{4} \\[6pt]
    y_G = m\, y_F = m \times 1 = m
    \end{cases}
    \]
    Ainsi :
    \[
    \mathbf{G\left(\frac{m}{4}\,;\, m\right)}.
    \]

    \item \textbf{Valeur de $m$ pour l'alignement des points $E$, $G$ et $B$ :}
    
    Les points $E(0\,;\,3)$, $B(1\,;\,0)$ et $G\left(\dfrac{m}{4}\,;\, m\right)$ sont alignés si et seulement si les vecteurs $\vect{EB}$ et $\vect{EG}$ sont colinéaires :
    \[
    E, G, B \text{ alignés} \iff \det(\vect{EG}, \vect{EB}) = 0.
    \]
    Déterminons les coordonnées de ces deux vecteurs :
    \[
    \vect{EB} = \begin{pmatrix} x_B - x_E \\ y_B - y_E \end{pmatrix} = \begin{pmatrix} 1 - 0 \\ 0 - 3 \end{pmatrix} = \begin{pmatrix} 1 \\ -3 \end{pmatrix},
    \]
    \[
    \vect{EG} = \begin{pmatrix} x_G - x_E \\ y_G - y_E \end{pmatrix} = \begin{pmatrix} \dfrac{m}{4} - 0 \\ m - 3 \end{pmatrix} = \begin{pmatrix} \dfrac{m}{4} \\ m - 3 \end{pmatrix}.
    \]
    Calculons le déterminant :
    \[
    \det(\vect{EG}, \vect{EB}) = \begin{vmatrix} \dfrac{m}{4} & 1 \\ m - 3 & -3 \end{vmatrix} = \frac{m}{4} \times (-3) - 1 \times (m - 3) = -\frac{3m}{4} - m + 3.
    \]
    Réduisons au même dénominateur :
    \[
    -\frac{3m}{4} - \frac{4m}{4} + 3 = -\frac{7m}{4} + 3.
    \]
    L'alignement équivaut donc à :
    \[
    -\frac{7m}{4} + 3 = 0 \iff \frac{7m}{4} = 3 \iff 7m = 12 \iff \mathbf{m = \frac{12}{7}}.
    \]
    \textbf{Conclusion :} Les points $E$, $G$ et $B$ sont alignés si et seulement si $\mathbf{m = \dfrac{12}{7}}$.
    
    \begin{center}
    \begin{tikzpicture}[scale=1.4, cap=round, >=Latex]
        % Sommets du parallélogramme
        \coordinate (A) at (0,0);
        \coordinate (B) at (2.5,0);
        \coordinate (D) at (0.8,1.4);
        \coordinate (C) at ($(B)+(D)-(A)$);
        
        % E tel que DE = 2 AD => AE = 3 AD
        \coordinate (E) at ($(A)!3!(D)$);
        
        % F tel que 3 FD + FC = 0 => F = D + 1/4 DC
        \coordinate (F) at ($(D)!0.25!(C)$);
        
        % Pour m = 12/7 : G = 12/7 F
        \coordinate (G) at ($(A)!{12/7}!(F)$);
        
        % Tracé du parallélogramme
        \draw[thick, fill=gray!8] (A) -- (B) -- (C) -- (D) -- cycle;
        
        % Repère (A; AB, AD)
        \draw[->, very thick, blue!80!black] (A) -- (B) node[midway, below] {\small $\vect{AB}$};
        \draw[->, very thick, blue!80!black] (A) -- (D) node[midway, left] {\small $\vect{AD}$};
        
        % Prolongement (AD) vers E
        \draw[dashed, gray!70] (D) -- (E);
        
        % Droite (AF) prolongée
        \draw[thick, blue!60!black] (A) -- ($(A)!2.1!(F)$) node[right] {\small $(AF)$};
        
        % Droite d'alignement (EB)
        \draw[very thick, dashed, red!80!black] ($(E)!-0.1!(B)$) -- ($(E)!1.1!(B)$) node[right] {\small $(EB)$};
        
        % Points et labels
        \fill (A) circle (1.5pt) node[below left] {$A$};
        \fill (B) circle (1.5pt) node[below right] {$B$};
        \fill (C) circle (1.5pt) node[above right] {$C$};
        \fill (D) circle (1.5pt) node[left] {$D$};
        \fill (E) circle (1.5pt) node[above left] {$E(0\,;\,3)$};
        \fill (F) circle (1.5pt) node[above] {$F\left(\frac{1}{4}\,;\,1\right)$};
        \fill (G) circle (1.5pt) node[below right] {$G\left(\frac{3}{7}\,;\,\frac{12}{7}\right)$};
    \end{tikzpicture}
    \end{center}
\end{enumerate}
\end{correctionbox}

\vspace{0.3cm}

% ====================================================
% EXERCICE 4
% ====================================================
\subsection*{Exercice 4 \hfill \normalfont\normalsize (4 points)}

Pour $m \in \R$, on considère la droite $d_m$ d'équation cartésienne :
\[
d_m : (m+2)x + (2m-1)y - m + 3 = 0.
\]
\textit{Les questions 1) et 2) sont à traiter de façon indépendante.}

\begin{enumerate}
    \item 
    \begin{enumerate}
        \item Montrer que $d_0$ et $d_1$ sont sécantes en un point $A$ dont on précisera les coordonnées.
        \item Montrer que $A$ est l'unique point appartenant à toutes les droites $d_m$ ($m \in \R$).
    \end{enumerate}
    \item 
    \begin{enumerate}
        \item Résoudre le système $\mathscr{S} : \begin{cases} x + 2y - 1 = 0 \\ 2x - y + 3 = 0 \end{cases}$
        \item Montrer que :
        \[
        (M(x\,;\,y) \in d_m \quad \forall m \in \R) \iff (x\,;\,y) \text{ est solution de } \mathscr{S}.
        \]
        \item Conclure.
    \end{enumerate}
\end{enumerate}

\begin{correctionbox}
\begin{enumerate}
    \item \textbf{Première approche : par analyse-synthèse}
    \begin{enumerate}
        \item \textbf{Intersection de $d_0$ et $d_1$ :}
        
        Pour $m = 0$ :
        \[
        d_0 : (0 + 2)x + (2 \times 0 - 1)y - 0 + 3 = 0 \iff 2x - y + 3 = 0.
        \]
        Pour $m = 1$ :
        \[
        d_1 : (1 + 2)x + (2 \times 1 - 1)y - 1 + 3 = 0 \iff 3x + y + 2 = 0.
        \]
        Un vecteur normal de $d_0$ est $\vect{n}_0(2\,;\,-1)$ et un vecteur normal de $d_1$ est $\vect{n}_1(3\,;\,1)$.
        Le déterminant de leurs coefficients directeurs vaut :
        \[
        \det(\vect{n}_0, \vect{n}_1) = \begin{vmatrix} 2 & 3 \\ -1 & 1 \end{vmatrix} = 2 \times 1 - (-1) \times 3 = 2 + 3 = 5 \neq 0.
        \]
        Les droites $d_0$ et $d_1$ ne sont donc pas parallèles : elles sont strictement sécantes en un unique point $A$.
        
        Déterminons les coordonnées de $A$ en résolvant le système :
        \[
        \begin{cases}
        2x - y + 3 = 0 \quad (L_1) \\
        3x + y + 2 = 0 \quad (L_2)
        \end{cases}
        \]
        En effectuant l'addition membre à membre $(L_1) + (L_2)$, les termes en $y$ s'annulent :
        \[
        (2x + 3x) + (-y + y) + (3 + 2) = 0 \iff 5x + 5 = 0 \iff 5x = -5 \iff x = -1.
        \]
        En remplaçant $x = -1$ dans $(L_1)$ :
        \[
        2(-1) - y + 3 = 0 \iff -2 - y + 3 = 0 \iff 1 - y = 0 \iff y = 1.
        \]
        Les droites $d_0$ et $d_1$ sont donc sécantes au point $\mathbf{A(-1\,;\,1)}$.

        \item \textbf{Unicité du point commun :}
        
        \begin{itemize}[leftmargin=*]
            \item \textbf{Condition nécessaire (Unicité) :}
            Si un point $M$ appartient à toutes les droites $d_m$ pour tout $m \in \R$, il appartient en particulier à $d_0$ et à $d_1$.
            Par suite, $M \in d_0 \cap d_1$. Or, d'après la question 1.(a), l'intersection $d_0 \cap d_1$ est réduite au seul point $A(-1\,;\,1)$.
            Par conséquent, $A$ est le \textbf{seul point candidat possible}.
            
            \item \textbf{Condition suffisante (Existence) :}
            Vérifions si les coordonnées de $A(-1\,;\,1)$ satisfont l'équation de $d_m$ quel que soit le paramètre réel $m$ :
            Pour tout $m \in \R$ :
            \begin{align*}
            (m+2)x_A + (2m-1)y_A - m + 3 &= (m+2)(-1) + (2m-1)(1) - m + 3 \\
            &= -m - 2 + 2m - 1 - m + 3 \\
            &= (-m + 2m - m) + (-2 - 1 + 3) \\
            &= 0m + 0 = 0.
            \end{align*}
            Cette égalité étant vraie pour tout $m \in \R$, le point $A(-1\,;\,1)$ appartient bien à $d_m$ pour tout réel $m$.
        \end{itemize}
        \textbf{Conclusion :} $A(-1\,;\,1)$ est l'\textbf{unique} point appartenant à toutes les droites $d_m$ pour tout $m \in \R$.
    \end{enumerate}

    \item \textbf{Deuxième approche : par identification polynomiale (méthode générale)}
    \begin{enumerate}
        \item \textbf{Résolution du système $\mathscr{S}$ :}
        \[
        \mathscr{S} : \begin{cases}
        x + 2y - 1 = 0 \quad (L_1) \\
        2x - y + 3 = 0 \quad (L_2)
        \end{cases}
        \]
        De l'équation $(L_1)$, on isole $x$ : $x = 1 - 2y$.
        
        En substituant dans $(L_2)$ :
        \[
        2(1 - 2y) - y + 3 = 0 \iff 2 - 4y - y + 3 = 0 \iff 5 - 5y = 0 \iff 5y = 5 \iff y = 1.
        \]
        On en déduit :
        \[
        x = 1 - 2(1) = 1 - 2 = -1.
        \]
        Le système $\mathscr{S}$ admet un unique couple solution :
        \[
        \mathbf{\mathcal{S} = \{(-1\,;\,1)\}}.
        \]

        \item \textbf{Équivalence logique :}
        
        Soit $M(x\,;\,y)$ un point du plan. Réorganisons l'expression de l'équation de $d_m$ en regroupant les termes selon les puissances de la variable $m$ :
        \begin{align*}
        (m+2)x + (2m-1)y - m + 3 &= mx + 2x + 2my - y - m + 3 \\
        &= m(x + 2y - 1) + (2x - y + 3).
        \end{align*}
        Ainsi :
        \[
        M(x\,;\,y) \in d_m \iff m(x + 2y - 1) + (2x - y + 3) = 0.
        \]
        Par suite, $M(x\,;\,y) \in d_m$ pour tout $m \in \R$ si et seulement si la fonction affine en $m$, définie par $P(m) = (x + 2y - 1)m + (2x - y + 3)$, est identiquement nulle sur $\R$.
        
        Or, un polynôme à coefficients réels de degré au plus 1 est nul pour tout $m \in \R$ si et seulement si tous ses coefficients sont nuls :
        \[
        \forall m \in \R, \quad m(x + 2y - 1) + (2x - y + 3) = 0 \iff \begin{cases} x + 2y - 1 = 0 \\ 2x - y + 3 = 0 \end{cases}
        \]
        ce qui signifie précisément que $(x\,;\,y)$ est solution du système $\mathscr{S}$.
        
        L'équivalence est ainsi rigoureusement démontrée :
        \[
        \mathbf{(M(x\,;\,y) \in d_m \quad \forall m \in \R) \iff (x\,;\,y) \text{ est solution de } \mathscr{S}}.
        \]

        \item \textbf{Conclusion :}
        
        Puisque le système $\mathscr{S}$ admet pour unique solution le couple $(-1\,;\,1)$, on retrouve le résultat :
        \[
        \mathbf{\text{Toutes les droites } d_m \text{ sont concourantes en l'unique point } A(-1\,;\,1)}.
        \]
    \end{enumerate}
\end{enumerate}
\end{correctionbox}

\begin{remarquebox}[Faisceau de droites et géométrie projective]
L'équation $m(x + 2y - 1) + (2x - y + 3) = 0$ s'interprète comme la combinaison linéaire des équations de deux droites fondamentales :
\[
\Delta_1 : x + 2y - 1 = 0 \quad \text{et} \quad \Delta_2 : 2x - y + 3 = 0.
\]
L'ensemble des droites $(d_m)_{m \in \R}$ forme ce que l'on appelle en mathématiques un \textbf{faisceau de droites de sommet $A$}.
Toutes les droites passant par $A$ appartiennent à ce faisceau, à l'exception de la droite limite obtenue lorsque $m \to \infty$, qui est précisément la droite $\Delta_1 : x + 2y - 1 = 0$.
\end{remarquebox}

\vspace{0.3cm}

% ====================================================
% EXERCICE 5
% ====================================================
\subsection*{Exercice 5 \hfill \normalfont\normalsize (3 points)}

On se place dans un repère orthonormé du plan.

On considère le point $A\left(\dfrac{1}{2}\,;\, \dfrac{9}{2}\right)$ et la droite $d$ d'équation :
\[
d : 3x - 4y + 4 = 0.
\]

\begin{enumerate}
    \item Soit $M(x\,;\,y)$ un point de la droite $d$. Exprimer $AM^2$ en fonction de $x$.
    \item En déduire la distance minimale entre $A$ et un point de la droite $d$.
\end{enumerate}

\begin{correctionbox}
\begin{enumerate}
    \item \textbf{Expression de $AM^2$ en fonction de $x$ :}
    
    Soit $M(x\,;\,y) \in d$. L'équation cartésienne de $d$ permet d'exprimer $y$ en fonction de $x$ :
    \[
    3x - 4y + 4 = 0 \iff 4y = 3x + 4 \iff y = \frac{3}{4}x + 1.
    \]
    Le repère étant orthonormé, le carré de la distance euclidienne $AM$ s'écrit :
    \[
    AM^2 = (x - x_A)^2 + (y - y_A)^2 = \left(x - \frac{1}{2}\right)^2 + \left(y - \frac{9}{2}\right)^2.
    \]
    En remplaçant $y$ par son expression $\dfrac{3}{4}x + 1$ :
    \[
    y - \frac{9}{2} = \left(\frac{3}{4}x + 1\right) - \frac{9}{2} = \frac{3}{4}x - \frac{7}{2}.
    \]
    Développons les deux carrés à l'aide des identités remarquables :
    \[
    \left(x - \frac{1}{2}\right)^2 = x^2 - 2 \times x \times \frac{1}{2} + \left(\frac{1}{2}\right)^2 = x^2 - x + \frac{1}{4},
    \]
    \[
    \left(\frac{3}{4}x - \frac{7}{2}\right)^2 = \left(\frac{3}{4}x\right)^2 - 2 \times \left(\frac{3}{4}x\right) \times \frac{7}{2} + \left(\frac{7}{2}\right)^2 = \frac{9}{16}x^2 - \frac{21}{4}x + \frac{49}{4}.
    \]
    Additionnons les deux expressions en regroupant les termes de même degré :
    \begin{align*}
    AM^2 &= \left(1 + \frac{9}{16}\right)x^2 + \left(-1 - \frac{21}{4}\right)x + \left(\frac{1}{4} + \frac{49}{4}\right) \\
    &= \left(\frac{16 + 9}{16}\right)x^2 + \left(-\frac{4 + 21}{4}\right)x + \frac{50}{4} \\
    &= \frac{25}{16}x^2 - \frac{25}{4}x + \frac{50}{4}.
    \end{align*}
    En simplifiant la constante $\dfrac{50}{4} = \dfrac{25}{2}$, on obtient l'expression demandée :
    \[
    \mathbf{AM^2 = \frac{25}{16}x^2 - \frac{25}{4}x + \frac{25}{2}} \quad \left(\text{ou } \frac{25}{16}x^2 - \frac{25}{4}x + \frac{50}{4}\right).
    \]

    \item \textbf{Distance minimale entre $A$ et la droite $d$ :}
    
    Posons $f(x) = \dfrac{25}{16}x^2 - \dfrac{25}{4}x + \dfrac{25}{2}$ pour $x \in \R$.
    
    $f$ est une fonction trinôme du second degré de la forme $f(x) = ax^2 + bx + c$, avec :
    \[
    a = \frac{25}{16}, \quad b = -\frac{25}{4}, \quad c = \frac{25}{2}.
    \]
    Puisque le coefficient dominant est strictement positif ($a = \dfrac{25}{16} > 0$), la fonction $f$ admet un minimum absolu sur $\R$, atteint en :
    \[
    x_0 = -\frac{b}{2a} = -\frac{-\dfrac{25}{4}}{2 \times \dfrac{25}{16}} = \frac{\dfrac{25}{4}}{\dfrac{25}{8}} = \frac{25}{4} \times \frac{8}{25} = \frac{8}{4} = \mathbf{2}.
    \]
    Calculons la valeur de ce minimum :
    \begin{align*}
    f(2) &= \frac{25}{16} \times 2^2 - \frac{25}{4} \times 2 + \frac{25}{2} \\
    &= \frac{25 \times 4}{16} - \frac{50}{4} + \frac{25}{2} \\
    &= \frac{25}{4} - \frac{25}{2} + \frac{25}{2} = \mathbf{\frac{25}{4}}.
    \end{align*}
    \textit{(Remarque par la forme canonique : $f(x) = \dfrac{25}{16}(x^2 - 4x) + \dfrac{25}{2} = \dfrac{25}{16}(x - 2)^2 - \dfrac{25}{4} + \dfrac{25}{2} = \dfrac{25}{16}(x - 2)^2 + \dfrac{25}{4} \ge \dfrac{25}{4}$, ce qui confirme sans calcul de dérivée le résultat)}.
    
    La fonction racine carrée $t \mapsto \sqrt{t}$ étant strictement croissante sur $[0\,;\, +\infty[$, la distance $AM = \sqrt{f(x)}$ atteint son minimum au même point $x_0 = 2$, et ce minimum vaut :
    \[
    d(A, d) = \min_{M \in d} AM = \sqrt{f(2)} = \sqrt{\frac{25}{4}} = \mathbf{\frac{5}{2} = 2{,}5}.
    \]
    \textbf{Conclusion :} La distance minimale entre le point $A$ et un point de la droite $d$ est égale à $\mathbf{\dfrac{5}{2}}$ (soit $2{,}5$).
\end{enumerate}
\end{correctionbox}

\begin{remarquebox}[Interprétation géométrique : Projeté orthogonal et formule de distance]
Le point $H$ de la droite $d$ réalisant la distance minimale est le \textbf{projeté orthogonal} de $A$ sur $d$.
\begin{itemize}[leftmargin=*,itemsep=1pt]
    \item Ses coordonnées sont : $x_H = 2$ et $y_H = \dfrac{3}{4}(2) + 1 = \dfrac{5}{2}$, soit $H\left(2\,;\, \dfrac{5}{2}\right)$.
    \item Le vecteur $\vect{AH}$ est :
    \[
    \vect{AH} = \begin{pmatrix} 2 - \dfrac{1}{2} \\ \dfrac{5}{2} - \dfrac{9}{2} \end{pmatrix} = \begin{pmatrix} \dfrac{3}{2} \\ -2 \end{pmatrix} = \frac{1}{2} \begin{pmatrix} 3 \\ -4 \end{pmatrix}.
    \]
    Le vecteur $\vect{n}\begin{pmatrix} 3 \\ -4 \end{pmatrix}$ est le vecteur normal canonique de la droite $d : 3x - 4y + 4 = 0$.
    On a donc $\vect{AH} = \dfrac{1}{2}\vect{n}$, ce qui prouve directement que $\vect{AH} \perp d$.
    \item La formule générale de la distance d'un point $A(x_A\,;\,y_A)$ à une droite d'équation $ax + by + c = 0$ donne immédiatement :
    \[
    d(A, d) = \frac{|ax_A + by_A + c|}{\sqrt{a^2 + b^2}} = \frac{\left|3 \times \dfrac{1}{2} - 4 \times \dfrac{9}{2} + 4\right|}{\sqrt{3^2 + (-4)^2}} = \frac{\left|\dfrac{3}{2} - 18 + 4\right|}{\sqrt{25}} = \frac{\left|-\dfrac{25}{2}\right|}{5} = \frac{\dfrac{25}{2}}{5} = \mathbf{\frac{5}{2}}.
    \]
    Ce résultat confirme avec une rigueur absolue la validité de nos calculs.
\end{itemize}

\begin{center}
\begin{tikzpicture}[scale=1.1, cap=round, >=Latex]
    % Grille et repère
    \draw[step=1cm, gray!20, very thin] (-1,-0.5) grid (4.5,5.5);
    \draw[->, gray!70] (-1,0) -- (4.8,0) node[right] {\small $x$};
    \draw[->, gray!70] (0,-0.5) -- (0,5.8) node[above] {\small $y$};
    
    % Droite d : 3x - 4y + 4 = 0 <=> y = 0.75 x + 1
    \draw[thick, blue!80!black] (-1, 0.25) -- (4.5, 4.375) node[above left] {\small $d : 3x - 4y + 4 = 0$};
    
    % Point A(0.5, 4.5)
    \coordinate (A) at (0.5, 4.5);
    % Point H(2, 2.5)
    \coordinate (H) at (2, 2.5);
    
    % Segment [AH] (distance minimale)
    \draw[very thick, red!80!black] (A) -- (H) node[midway, right] {\small $d(A,d) = \frac{5}{2}$};
    
    % Marqueur d'angle droit en H
    % Vecteur u unitaire le long de d : (0.8, 0.6)
    % Vecteur n unitaire vers A : (-0.6, 0.8)
    \draw[thick, gray!80] ($(H) + 0.3*(0.8, 0.6)$) -- ($(H) + 0.3*(0.8, 0.6) + 0.3*(-0.6, 0.8)$) -- ($(H) + 0.3*(-0.6, 0.8)$);
    
    \fill (A) circle (1.8pt) node[above] {$A\left(\frac{1}{2}\,;\,\frac{9}{2}\right)$};
    \fill (H) circle (1.8pt) node[below right] {$H\left(2\,;\,\frac{5}{2}\right)$};
\end{tikzpicture}
\end{center}
\end{remarquebox}

\end{document}
